\documentstyle[% 


righttag% 


,11pt% 


,amssymb% 


,russian%               remove in Eng. 


,izvru%                 Set izven% in Eng. 


% ,izvru-ks             for brief communications 


]{amsart} 


 


% 


 


\newcommand{\tts}[1]{} 


 


\theoremstyle{definition} 


\theorembodyfont{\rm} 


\newtheorem{cornonum}{\hskip\parindent Следствие} 


\renewcommand{\thecornonum}{} 


%\hoffset=-15mm%                  For Eng. 


 


\setcounter{page}{81} 


\god{2000} 


\nomer{7~(458)} %                        Remove in Eng. 


%\nomer{Vol.\,44, No.\,7}%               Set in Eng. 


%\str{\pageref{begin}--\pageref{end}}%  Set in Eng. 


\udk{517.987} 


 


\author{\it В.И.\,Чилин, И.Г.\,Ганиев} 


% NB:   ^^ remove in Eng. 


\title{Индивидуальная эргодическая теорема для сжатий \\ 


в решетке $L_p(\wh\nabla,\wh\mu)$ Банаха--Канторовича} 


 


\date{} 


% \pagestyle{myheadings}%         Set in Eng. 


\pagestyle{plain} %              Remove in Eng. 


\begin{document} 


% \thispagestyle{plain}%          Set in Eng. 


\maketitle 


\label{begin} 


 


В работе получен аналог известной эргодической теоремы для 


положительных сжатий решеток 


$L_p(\wh\nabla,\wh\mu)$ Банаха--Канторовича, построенных 


по мере $\wh\mu$ со значениями в кольце измеримых функций. 


 


Пусть $(\Omega,\Sigma,\lambda)$ -- пространство с конечной мерой, 


$L_0(\Omega)$ --- кольцо 


измеримых функций на $\Omega$ (равные $\lambda$-почти всюду функции 


отождествляются), $\nabla(\Omega)$ --- булева алгебра всех идемпотентов 


в $L_0(\Omega)$, 


$\wh\nabla$ --- произвольная полная булева алгебра, содержащая 


$\nabla(\Omega)$ как 


правильную подалгебру, $\wh\mu:\wh\nabla\to L_0(\Omega)$ --- строго 


положительная 


$L_0(\Omega)$-значная мера на $\wh\nabla$ [1], причем $\wh\mu(g\wh 


e)=g\wh\mu(\wh e)$ 


для всех $\wh e\in\wh\nabla$, $g\in\nabla(\Omega)$. 


 


Обозначим через $X(\wh\nabla)$ экстремальный вполне несвязный компакт, 


отвечающий булевой алгебре $\wh\nabla$, а через $L_0(\wh\nabla)$ --- 


алгебру всех непрерывных 


действительных функций, заданных на $X(\wh\nabla)$, принимающих 


значения $\pm\infty$ на нигде неплотных множествах из $X(\wh\nabla)$ 


[2]. Ясно, что $L_0(\Omega)$ является подалгеброй в $L_0(\wh\nabla)$. 


 


Следуя известной схеме построения $L_p$-пространств, определим 


пространство $L_p(\wh\nabla,\wh\mu)$ 


всех функций из $L_0(\wh\nabla)$, интегрируемых с 


$p$-й степенью по $L_0(\Omega)$-значной мере $\wh\mu$ (см., напр., 


[3]). 


 


Известно [4], что $L_p(\wh\nabla,\wh\mu)$ 


является решеткой Банаха--Канторовича 


относительно соответствующей $L_0(\Omega)$-значной нормы $|\cdot|_p$, 


при этом $L_p(\wh\nabla,\wh\mu)$ 


--- модуль над $L_0(\Omega)$. 


 


Естественно, что пространства 


$L_p(\wh\nabla,\wh\mu)$ должны обладать многими свойствами, 


которые аналогичны соответствующим свойствам классических 


функциональных $L_p$-пространств, построенных по числовым мерам. 


Доказательства таких свойств могут осуществляться с использованием 


одного из следующих трех методов. 


 


1) Последовательное повторение всех шагов известных доказательств 


для классических $L_p$-пространств, учитывающее особенности 


$L_0(\Omega)$-значных мер. 


 


2) Применение булевозначного анализа, дающего возможность в 


соответствующей модели теории множеств сводить $L_0(\Omega)$-модули 


$L_p(\wh\nabla,\wh\mu)$ 


к обычным банаховым $L_p$-пространствам. 


 


3) Представление $L_p(\wh\nabla,\wh\mu)$ 


в виде измеримого расслоения классических 


$L_p$-пространств, опирающееся на существование соответствующих 


лифтингов. 


 


Первый метод, разумеется, неэффективен, поскольку он вынуждает 


повторять известные этапы доказательств, модифицируя их для 


$L_0(\Omega)$-значных мер. Второй метод связан с привлечением 


достаточно 


трудоемкого аппарата булевозначного анализа, и его реализация потребует 


огромной подготовительной работы, связанной с установлением 


взаимосвязей для изучаемых объектов в обычной и булевозначной моделях 


теории множеств. 


 


Наиболее ``экономным'' путем для исследования свойств пространств 


$L_p(\wh\nabla,\wh\mu)$ 


с нашей точки зрения является использование третьего метода. 


Это связано с тем, что уже имеется достаточно хорошо разработанная 


теория измеримых расслоений банаховых пространств (см., 


напр., [5]). Поэтому ее эффективное использование с привлечением 


уже известных результатов из общей теории банаховых пространств 


дает возможность получения различных свойств и для пространств 


Банаха--Канторовича. 


 


В данной работе этот метод иллюстрируется на примере доказательства 


эргодической теоремы для положительных сжатий решеток 


$L_p(\wh\nabla,\wh\mu)$ 


Банаха--Канторовича. 


 


Пусть $\wh\nabla$, $\wh\mu$ те же, что и выше. 


Рассмотрим измеримое расслоение булевых алгебр $\nabla_\omega$ со 


строго 


положительной числовой мерой $\mu_\omega$, $\omega\in\Omega$, 


соответствующее 


$(\wh\nabla,\wh\mu)$ [6], [7]. 


Пусть $L_0(\wh\nabla)=L_0(\wh\nabla,\wh\mu)$ --- топологическая 


векторная решетка 


функций с $L_0(\Omega)$-значной метрикой $\wh\rho(\wh f,\wh 


g)=\int\limits_{\wh\nabla}\frac{|\wh f-\wh g|}{1+|\wh f-\wh g|}d\wh\mu$ 


и $L_0(\nabla_\omega,\mu_\omega)$ --- измеримое расслоение 


топологических векторных 


решеток измеримых функций, соответствующее 


$L_0(\wh\nabla,\wh\mu)$ с метрикой 


$\rho_\omega(f(\omega),g(\omega))=\int\limits_{\nabla_\omega} 


\frac{|f(\omega)-g(\omega)|}{1+|f(\omega)-g(\omega)|}d\mu_\omega$ 


(см. [8]); здесь и в дальнейшем $\wh f$ означает класс 


эквивалентности из 


$L_0(\wh\nabla,\wh\mu)$, 


содержащий представитель $f(\omega)\in L_0(\nabla_\omega,\mu_\omega)$, 


$\omega\in\Omega$. 


 


\begin{thm} 


Для последовательности $\{\wh f_n\}$ из 


$L_0(\wh\nabla,\wh\mu)$ точная 


верхняя $($соответственно нижняя$)$ грань $\sup\limits_{n\ge1}\wh f_n$ 


$($соответственно 


$\inf\limits_{n\ge1}\wh f_n)$ существует в 


$L_0(\wh\nabla,\wh\mu)$ 


в том и только том случае, 


когда существует $\sup\limits_{n\ge1}f_n(\omega)$ $($соответственно 


$\inf\limits_{n\ge1}f_n(\omega))$ в 


$L_0(\nabla_\omega,\mu_\omega)$ для почти всех $($п.\,в.$)$


$\omega\in\Omega$. В этом 


случае 


$(\sup\limits_{n\ge1}\wh f_n)(\omega)= 


\sup\limits_{n\ge1}f_n(\omega)$ 


$\big((\inf\limits_{n\ge1}\wh f_n)(\omega)= 


\inf\limits_{n\ge1}f_n(\omega)\big)$ для п.\,в. $\omega\in\Omega$. 


\end{thm} 


 


Из этой теоремы вытекает 


 


\begin{cornonum} 


Если последовательность $\wh f_n$ элементов из решетки 


$L_0(\wh\nabla,\wh\mu)$\; 


$(o)$-сходится к $\wh f\in L_0(\wh\nabla,\wh\mu)$ 


(запись: $\wh f_n@>(o)>>\wh f$), 


то $f_n(\omega)@>(o)>> f(\omega)$ в 


$L_0(\nabla_\omega,\mu_\omega)$ для п.\,в. $\omega\in\Omega$. Обратно, 


если 


$f_n(\omega)@>(o)>> g(\omega)$ для некоторого 


$g(\omega)\in L_0(\nabla_\omega,\mu_\omega)$ для п.\,в. 


$\omega\in\Omega$, 


то $\wh g\in L_0(\wh\nabla,\wh\mu)$ и $\wh f_n@>(o)>>\wh g$ в 


$L_0(\wh\nabla,\wh\mu)$. 


\end{cornonum} 


 


Пусть 


$T:L_p(\wh\nabla,\wh\mu)\to L_p(\wh\nabla,\wh\mu)$ 


--- линейное отображение. Как обычно, 


$T$ называют положительным, если $T\wh f\ge0$ для всех $\wh f\ge0$. 


Предположим, что $T-L_0(\Omega)$ --- ограниченное отображение, т.\,е. 


существует такое $k\in L_0(\Omega)$, что $|T\wh f|_p\le k|\wh f|_p$ при 


всех 


$\wh f\in L_0(\wh\nabla,\wh\mu)$. 


В этом случае определен элемент 


$|T|=\sup\limits_{|\wh f|_p\le\bold 1}|T\wh f|_p$, 


который является элементом $L_0(\Omega)$, где $\bold 1$ --- единичный 


элемент 


$L_0(\Omega)$. Если $|T|\le\bold 1$, то отображение называют сжатием. 


\medskip 


 


{\bf Пример.} Пусть $(\Omega,\wh\nabla,\lambda)$ --- измеримое 


пространство с конечной 


числовой мерой $\lambda$. Под $\wh\nabla$ подразумеваем полную булеву 


алгебру 


всех классов эквивалентности равных $\lambda$-п.\,в. измеримых множеств 


из $\Omega$. Пусть $\wh\nabla_0$ --- правильная булева подалгебра в 


$\wh\nabla$, 


$\lambda_0$ --- сужение $\lambda$ на $\wh\nabla_0$, 


$E(\cdot\mid\wh\nabla_0)$ --- условное математическое 


ожидание из $L_1(\Omega,\wh\nabla,\lambda)$ на 


$L_1(\Omega,\wh\nabla_0,\lambda_0)$. Ясно, что $\wh\mu(\wh e)=E(\wh 


e\mid\wh\nabla_0)$ 


является строго положительной 


$L_0(\Omega,\wh\nabla_0,\lambda_0)$-значной мерой на $\wh\nabla$. 


 


Пусть $\wh\nabla_1$ --- другая произвольная правильная булева 


подалгебра в 


$\wh\nabla$ и $\wh\nabla_1\supset\wh\nabla_0$. Обозначим через 


$\wh\mu_1$ сужение $\wh\mu$ на $\wh\nabla_1$. 


Согласно теореме 4.2.9 ([4], с.\, 182) существует оператор условного 


математического ожидания 


$T:L_1(\wh\nabla,\wh\mu)\to L_1(\wh\nabla_1,\wh\mu_1)$, который 


является положительным и отображает 


$L_p(\wh\nabla,\wh\mu)$ 


в $L_p(\wh\nabla,\wh\mu)$ для всех $p\ge1$, при этом $|T\wh f|_p\le|\wh 


f|_p$ для любого 


$\wh f\in L_p(\wh\nabla,\wh\mu)$ и $T\bold 1=\bold 1$. 


 


\begin{thm} 


Пусть $T:L_p(\wh\nabla,\wh\mu)\to L_p(\wh\nabla,\wh\mu)$ 


--- положительное линейное 


сжатие и $T\bold 1{\le}\bold 1$. Тогда для каждого $\omega\in\Omega$ 


существует 


такое положительное сжатие 


$T_\omega:L_p(\nabla_\omega,\mu_\omega)\to 


L_p(\nabla_\omega,\mu_\omega)$, 


что $T_\omega f(\omega)=(T\wh f)(\omega)$ для любого 


$\wh f\in L_p(\wh\nabla,\wh\mu)$ и $\lambda$-п.\,в. $\omega$. 


\end{thm} 


 


Следующая теорема является аналогом известной эргодической 


теоремы для положительных сжатий решеток 


$L_p(\wh\nabla,\wh\mu)$ Банаха--Канторовича, 


построенных по $L_0(\Omega)$-значным мерам. 


 


\begin{thm} 


Пусть $p>1$, $q>1$, $\frac{1}{p}+\frac{1}{q}=1$, 


$T:L_p(\wh\nabla,\wh\mu)\to L_p(\wh\nabla,\wh\mu)$ --- 


положительное линейное сжатие, для которого $T\bold 1\le \bold 1$. 


Тогда 


для каждого $\wh f\in L_p(\wh\nabla,\wh\mu)$ 


\begin{itemize} 


\item[1)] последовательность $s_n(|\wh 


f|)=\frac{1}{n}\sum\limits^{n-1}_{i=0}T^i(|\wh f|)$ ограничена в 


$L_p(\wh\nabla,\wh\mu)$ и $|\sup\limits_{n\ge1}s_n(|\wh f|)|_p\le q|\wh 


f|_p;$ 


\item[2)] существует такое 


$\wt f\in L_p(\wh\nabla,\wh\mu)$, что последовательность 


$s_n(\wh f)$\; $(o)$-сходится к $\wt f$ в 


$L_p(\wh\nabla,\wh\mu)$, в частности, 


$s_n(\wh f)$\; $\wh\mu$-п.\,в. сходится к $\wt f$. 


\end{itemize} 


\end{thm} 


 


\begin{pf} 


Пусть $\wh f\in L_p(\wh\nabla,\wh\mu)$ и $s_n(\wh 


f)=\frac{1}{n}\sum\limits^{n-1}_{i=0}T^i\wh f$. 


Согласно теореме 2\; $s_n(\wh 


f)(\omega)=\frac{1}{n}\sum\limits^{n-1}_{i=0} T^i_\omega 


f(\omega)=s_n(f(\omega))$ для п.\,в. $\omega$. 


Из известных теорем 5.25 и 5.26 ([9], с.\,189--190) следует, 


что $\|\sup\limits_{n\ge1}s_n(f(\omega))\|_p\le q\|f(\omega)\|_p$ и 


существует такое 


$g(\omega)\in L_p(\nabla_\omega,\mu_\omega)$, что 


$s_n(f(\omega))@>(o)>>g(\omega)$ в $L_p(\nabla_\omega,\mu_\omega)$ 


для п.\,в. $\omega\in\Omega$. 


 


В силу теоремы 1 в $L_p(\wh\nabla,\wh\mu)$ существует 


$\sup\limits_{n\ge1}s_n(|\wh f|)$ и 


$|\sup\limits_{n\ge1}s_n(|\wh f|)|_p\le q|\wh f|_p$. Из следствия к 


теореме 1 


вытекает, что $s_n(\wh f)=\wh{s_n(f(\omega))}@>(o)>>\wt 


f=\wh{g(\omega)}$ в $L_0(\wh\nabla,\wh\mu)$. 


Поскольку последовательность $s_n(\wh f)$ ограничена в 


$L_p(\wh\nabla,\wh\mu)$, 


то $s_n(\wh f)@>(o)>>\wt f$ в $L_p(\wh\nabla,\wh\mu)$. 


Так же, как и в [10], доказывается, что из $(o)$-сходимости в 


$L_0(\wh\nabla,\wh\mu)$ 


следует сходимость $\wh\mu$-п.\,в. Поэтому $s_n(\wh f)\to\wt f$\; 


$\wh\mu$-п.\,в. 


\end{pf} 


 


\begin{thebibliography}{99} 


 


\bibitem{1} 


Сарымсаков Т.А., Рубштейн Б.А., Чилин В.И. {\em Полные тензорные 


произведения топологических полуполей} // ДАН СССР. -- 1974. -- 


Т.\,216. -- \No\,6. -- С.\,1226--1228. 


\bibitem{2} 


Сарымсаков Т.А., Аюпов Ш.А., Хаджиев Дж., Чилин В.И. {\em Упорядоченные 


алгебры}. -- Ташкент: Фан, 1983. -- 304 с. 


\bibitem{3} 


Сарымсаков Т.А. {\em Топологические полуполя и их приложения}. -- 


Ташкент: Фан, 1989. -- 192~с. 


\bibitem{4} 


Кусраев А.Г. {\em Векторная двойственность и ее приложения}. -- 


Новосибирск: Наука, 1985. -- 256 с. 


\bibitem{5} 


Гутман А.Е. {\em Банаховы расслоения в теории решеточно нормированных 


пространств. Линейные операторы, согласованные с порядком} // 


Тр. РАН Сиб. отд. ин-та матем. -- 1995. -- Т.\,29. -- 


C.\,63--211. 


\bibitem{6} 


Ганиев И.Г. {\em О векторных мерах со значениями в пространствах 


Банаха--Канторовича} // Изв. вузов. Математика. -- 1999. 


-- \No\,4. -- С.\,65--67. 


\bibitem{7} 


Ганиев И.Г. {\em Измеримые булевы расслоения} // Тез. III Сиб. 


конгр. по прикл. и индустр. матем. -- Новосибирск, 1998. -- 


Ч.\,1. -- С.\,65. 


\bibitem{8} 


Ганиев И.Г. {\em Измеримые расслоения метризуемых топологических 


векторных решеток} // ДАН РУз. -- 1999. -- \No\,3. -- С.\,8--11. 


\bibitem{9} 


Krendel U. {\em Ergodic Theorems}. -- Walter de Gruywer. -- 


Berlin--N.\,Y., 1985. -- 357 p. 


\bibitem{10} 


Сарымсаков Т.А. {\em Полуполя и теория вероятностей}. -- Ташкент: 


Фан, 1980. -- 94 с. 


 


\end{thebibliography} 


 


\vskip 2cm 


 


\post{\it Ташкентский государственный университет}{\it Поступила} 


\post{\it Ташкентский институт инженеров}{19.11.1999} 


\post{\it железнодорожного транспорта}{} 


 


\label{end} 


\end{document} 


